\documentclass[../numerical_methods.tex]{subfiles}
\begin{document}
\section{Aula 07 - 24 de Março, 2026}
\subsection{Motivações}
\begin{itemize}
	\item Organização de Código;
	\item Funções e Módulos.
\end{itemize}
\subsection{Organização de Código, Funções e Módulos}
Vamos aprender, hoje, a sair dos scripts gigantes e começar a organizar nosso código matemático em blocos lógicos reutilizáveis: as \textit{funções}, os \textit{módulos} e as \textit{subrotinas}.

As funções já são conhecidas por nós no Python -- elas recebem entradas (parâmetros), realizam os cálculos e retornam um valor.
\begin{listing}
	\caption{Calculando a Distância entre Dois Pontos}
	\begin{minted}[bgcolor = DarkBackground, linenos=true, breaklines]{Python}
  import math 
  def calcular_distancia(x1, y1, x2, y2):
      """ Calcula a distancia euclidiana entre P1=(x1, y1) e P2=(x2, y2)."""
      distancia = math.sqrt((x2-x1)**2+(y2-y1)**2)
      return distancia

  # Teste
  d = calcular_distancia(1.0, 2.0, 4.0, 6.0)
  print(f"A distância entre os pontos 1 e 2 é: {d}")
  \end{minted}
\end{listing}
Diferente da matemática tradicional, uma função em Python pode retornar vários valores de uma vez usando tuplas.
\begin{minted}[bgcolor = DarkBackground, linenos=true, breaklines]{Python}
## Função retornando múltiplos valores 
def analisar_dados(lista_valores):
    """Função Matemática: retorna média, máximo E mínimo"""
    if len(lista_valores)==0:
        return 0,0,0

    media = sum(lista_valores)/len(lista_valores)
    return media, max(lista_valores), min(lista_valores)
\end{minted}

Podemos criar, também, \textbf{subrotinas}, que são funções responsáeis por retornar valores não matemáticos, mas sim executar uma ação visual, como um relatório.
\begin{listing}[H]
	\caption{Subrotina: formatar o código de um boletim.}
	\begin{minted}[bgcolor = DarkBackground, linenos=true, breaklines]{Python}
def imprimir_boletim(nome_turma, notas):
    """Subrotina visual: chama a função matemática e formata a tela"""
    media, nota_max, nota_min = analisar_dados(notas)
    printf("-"*30)
    print(f"Turma: {note_turma}")
    print(f"Média: {media:.2f} | Max: {nota_max:.2f} |
    Min:{nota_min:.2f}")
    print("-"*30)
    # o .2f arredonda o valor para duas casas decimais!

# O uso fica muito mais limpo: 
notas_turma = [4.5, 6.0, 3.2, 8.5, 5.0, 7.5]
imprimir_boletim("Cálculo numérico: ", notas_turma)
\end{minted}
\end{listing}

Se uma regra matemática exige uma condição, nosso código deve ser capaz de impôr essa condição antes de tentar fazer a conta; um dos exemplos clássicos disso é a regra de Cramer -- vamos criar uma função para resolver um sistema 2x2 do tipo
\[
	\left\{\begin{array}{ll}
		ax +by = e \\
		cx + dy = f
	\end{array}\right..
\]
Pela álgebra linear, sabemos que o determinante da matriz dos coeficientes deve ser não nulo para que esse sistema seja resolvível, então vamos usar o comando \texttt{raise ValueError} para blindar nossa função e avisar do erro!
\begin{listing}[H]
	\caption{Calculo de Determinante Sensível a Erro}
	\begin{minted}[bgcolor = DarkBackground, linenos=true, breaklines]{Python}
def resolver_sistema(a, b, c, d, e, f):
    """
    Resolve o sistema linear: 
    ax + by = e 
    cx + dy = f
    """
    # 1. Calculamos o determinante 
    det_principal = (a*d)-(b*c)
    
    # 2. A Barreira de Rigor Matemático!
    if det_principal==0:
       raise ValueError("Erro Matemático: o determinante é zero. Sistema impossível ou indeterminado")
    # 3. Se passou pela barreiar, o cálculo é seguro: 
    x = ((e*d)-(b*f))/det_principal 
    y = ((a*f)-(e*c))/det_principal
    
    return x, y

# Teste 1: um sistema com solução (x=2, y=3)
x_res, y_res = resolver_sistema(1,2,3,4,8,18)
print(f"Solução do sistema: x={x_res}, y={y_res}")

# Teste 2: um sistema sem solução única
resolver_sistema(1,1,2,2,5,10)
 
 \end{minted}
\end{listing}

Agora vamos aprender a organizar os arquivos como um todo criando nossas próprias bibliotecas: simularemos a criação de arquivo separado \texttt{algebra.py}, que pode ser feito com o comando \texttt{\%\%writefile nome\_arquivo.py} no Colab, logo na primeira linha da célula; depois disso, ele aparecerá na pasta lateral do Colab.
\begin{listing}[H]
	\caption{Criação de módulo no Colab}
	\begin{minted}[bgcolor = DarkBackground, linenos=true, breaklines]{Python}
 %%writefile algebra.py 

 def multiplicar_matrizes(A, B):
    """ Multiplica a matriz A (m x n) pela B (n x p). """
    linhas_A=len(A)
    colunas_A = len(A[0])
    linhas_B=len(B)
    colunas_B = len(B[0])

    if colunas_A != linhas_B:
        raise ValueError("Erro: colunas devem ser iguais às linhas de B.")
 
 # Equivalente a: 
 # import numpy as np 
 # C = npzeros((linhas_A, colunas_B)) 

 C = []
 for i in range(linhas_A):
    C.append([0.0]*colunas_B)

 for i in range(linhas_A):
    for j in range(colunas_B):
        soma = 0
        for k in range(colunas_A):
            soma+=A[i][k]*B[k][j]
        C[i][j]=soma
  return C

  def imprimir_matriz(nome, matriz):
      """Imprime a matriz formatada"""
      print(f"Matriz {nome}: ")
      for linha in matriz: 
          linha_formatada = [f"{elemento:>5.1f}" for elemento in linha]
          print(" [" + "".join(linha_formatada)+"]")
      prit("-"*15) 
 \end{minted}
\end{listing}

Abaixo, vamos agir como se estivéssemos em outro arquivo e importaremos a biblioteca que criamos!
\begin{listing}[H]

	\begin{minted}[bgcolor = DarkBackground, linenos=true, breaklines]{Python}
import algebra as alg 

matriz_A = [[1.0, 2.0], [3.0, 4.0]]
matriz_B = [[1.0, 2.0], [3.0, 4.0]]

alg.imprimir_matriz("A", matriz_A)
alg.imprimir_matriz("B", matriz_B)

resultado = alg.multiplicar_matrizes(matriz_A, matriz_B)
alg.imprimir_matriz("C {resultado}", resultado)
\end{minted}
\end{listing}

\subsection{Aplicação: cifra de Hill}
\textbf{\underline{O Problema:}} queremos enviar a mensagem "BOM DIA " com um espaço no final pela rede, mas não queremos que os hackers leiam.
Para isso, vamos transformar letras em números (Espaço)=0, A=1, B=2, etc. Se agruparmos esses números de dois em dois nas colunas, formamos a nossa \textbf{matriz da mensagem M}:
\[
	M = \begin{bmatrix}
		2  & 13 & 4  & 1 \\
		15 & 0  & 90 & 0
	\end{bmatrix}
\]

\textbf{\underline{A Senha:}} o remetente e o destinatário combinam uma senha secreta, que é uma matriz 2x2
\[
	K = \begin{bmatrix}
		2 & 1 \\
		1 & 1
	\end{bmatrix},
\]
então para trancar a mensagem, basta usar o produto da chave pela mensagem: \(C = K \times M\)
\begin{tcolorbox}[
		skin=enhanced,
		title=Observação,
		fonttitle=\bfseries,
		colframe=black,
		colbacktitle=cyan!75!white,
		colback=cyan!15,
		colbacklower=black,
		coltitle=black,
		drop fuzzy shadow,
		%drop large lifted shadow
	]
	Como nosso alfabeto do exercício tem 27 caracteres, pegamos o resultado e dividimos por 27, guardando apenas o resto da divisão; isto garante que o número sempre volte a ser uma letra!
\end{tcolorbox}

\textbf{\underline{Como Destrancar:}} não existe ``divisão'' de matrizes, então é necessário fazer o produto pela inversa.
\begin{listing}[H]
	\begin{minted}[bgcolor = DarkBackground, linenos=true, breaklines]{Python}
%%writefile motor_cripto.py 
"""
====================================== 
Módulo: motdor de criptografia 
Objetivo: conter as funções matemática e de tradução 
======================================
"""

ALFABETO = " ABCDEFGHIJKLMNOPQRSTUVWXYZ"

def texto_para_matriz(texto):
    """Tranforma um texto par em uma matriz com 2 linhas"""

    if len(texto)%2!=0:
        texto += " "

    colunas = len(texto)//2 
    matriz = [[0]*colunas, [0]*colunas]

    for i in range (colunas):
        matriz[0][i]=ALFABETO.index(texto[i*2])
        matriz[1][i]=ALFABETO.index(texto[i*2+1])

    return matriz

def operar_matriz_modular(matriz_A, matriz_B):
    """
    Tem que multiplicar A por B e obter o resto da divisão por 27.
    """
    linhas_A=len(matriz_A)
    linhas_B=len(matriz_B)
    colunas_A=len(matriz_A[0])
    colunas_B=len(matriz_B[0])

    # Estrutura de criação de matriz nula com python puro.
    matriz_resultado = []
    for i in range(linhas_A):
        matriz_resultado.append([0]*colunas_B)
\end{minted}
\end{listing}

\begin{minted}[bgcolor = DarkBackground, linenos=true, breaklines]{Python}

    # Produto matricial 
    for i in range(linhas_A):
        for j in range(colunas_B):
            soma=0
            for k in range(linhas_B):
                soma+=matriz_A[i][k]*matriz_B[k][j] 

            matriz_resultado[i][j]=soma%27
    return matriz_resultado

  # Resta converter de volta para texto: 
  def matriz_para_texto(matriz):
      """
      Lê uma matriz e converte seus números em letras.
      """
      linhas = len(matriz)
      colunas= len(matriz[0])

      texto_final=""
      for j in range(colunas):
          for i in range(linhas):
              numero=matriz[i][j]
              texto_final+=ALFABETO[numero]
      return texto_final
\end{minted}
Tendo o motor cripto montado, temos que montar o script principal, nomeado \textbf{app espião} (shhhh); para tanto, basta fazer

\begin{listing}[H]
	\begin{minted}[bgcolor = DarkBackground, linenos=true, breaklines]{Python}
 # ============================== 
 # NASCE UM ESPIÃO
 # ============================== 

 import motor_cripto as cripto

 # 1. As chaves matemáticas 
  matriz_chave = [
      [2,1],
      [1,1]
  ]
  matriz_inversa = [
  [1,26],
  [26, 1]
  ]

  mensagem_original = "BOM DIA"

  print("="*30)
  print("SISTEMA DE COMUNICAÇÃO SEGURA INICIADO")
  print("="*30)
  print(f"Código secreto que precisamos transmitir: '{mensagem_original}'\n")

  # ==================
  # Fase 1: REMETENTE
  # ==================
  print(">> FASE 1: criptografando segredos...")
  matriz_M = cripto.texto_para_matriz(mensagem_original)

  matriz_C = cripto.operar_matrizes(matriz_chave, matriz_M)
  mensagem_hacker=cripto.matriz_para_texto(matriz_C)

  prit("Segredos Indecifráveis:")
  for linha in matriz_C: 
      print(linha)
  print(f"TEXTO INTERCEPTADO NA REDE: '{mensagem_hacker}'\n")
\end{minted}
\end{listing}


\begin{minted}[bgcolor = DarkBackground, linenos=true, breaklines]{Python}
  # =======================
  # Fase 2: O DESTINATÁRIO 
  # =======================

  print(">> FASE 2: desvendando os mistérios.")

  matriz_C_recebida = cripto.texto_para_matriz(mensagem_hacker)
  # Pulo do gato: multiplicar a matriz bagunçada pela inversa da chave!
  matriz_M_recuperada = cripto.operar_matrizes(matriz_inversa, matriz_C_recebida)
  mensagem_salva = cripto.matriz_para_texto(matriz_M_recuperada)

  print("Matriz Recuperada:")
  for linha in matriz_M_recuperada:
      print(linha)
  print(f"Mensagem desvendada com sucesso: '{mensagem_salva}'")
 \end{minted}
\end{document}
